% Pursuing Stacks (cotes 134-2 to 134-8): archive cote 134-3, site page 9 (PDF page 10), typescript page 9.
% typescript. Draft compared: ps2.tex:372-433.
% AI-transcribed by Claude Opus 5.5 (claude-opus-5-5, adaptive thinking, effort high), Message Batches API, batch ps1 (msgbatch_01Rv5hCaNTtUt4JHx1yqoHPL), 2026-11-12; not yet checked by eye.
\scanpage{134-3/9}
\margin{9}
\mspage{9}

These reflections suggest that there should be a rather impressive
bunch of algebraic structures,  each giving rise to a model category
for \inserted{the} usual homotopy category,  or in any case yielding this category by
localization with respect to a suitable notion of "weak equivalences".
T\unsure{\over{h}{h}}e"bunch"is all the more impressive,  if we remember that the notion
of a stack  (dropping now the qualification $n$ namely assuming $n=\infty$ )  is
not really a uniquely defined one,  but depends on the choice of a
"\uline{coherator}",  namely (mainly) a category $C$ satisfying certain require-%
ments,  which can be met in a vast variety of ways\hand{,} presumably.  T\unsure{\over{h}{h}}e con-%
struction of coherators is achieved in terms of universal algebra,
which seems here the undispensable Ariane thread not to get lost in
overwhelming messyness.  T\unsure{\over{h}{h}}e natural question which arises\inserted{here} (and which I
\margin{\unsure{$\lambda$} (10)}\uline{do not feel though like pursuing})\hand{\unsure{//}} is to give\inserted{in terms of universal algebra} some kind of caracterizati-%
on,  among all "algebraic" structures,  of those which give rise in
some specified ways\inserted[typed]{(including the known \unsure{ones|cases})} to a category of models say for usual homotopy theo-%
ry.

1\unsure{\over{8}{7}}.  When referring \deleted{\illegible{1}}\inserted{(section 15)} to the notion of a stack as a unifying concept
for homotopical algebra and non commutative cohomological algebra,  I
forgot to mention one signfficant \deleted{\illegible{1}}\inserted{observation} of Artin-Mazur along those li-%
nes\inserted{(\unsure{namely|messy}, unification)},  namely that (for ordinary homotopy types say) weak equivalences
(namely maps inducing isomorphisms for all homotopy groups)  can be ca-%
racterized as being those which induce isomorphisms on cohomology groups
\xout{\illegible{1}} of the spaces considered not only for constant coefficients,  but
also for \deleted{\xout{\illegible{1}}} arbitrary twisted coefficients on the target space,
including also the non commutative H$^{0}$ and H$^{1}$ for twisted \xout{\illegible{1}} (non commu%
% NOTE: the end of this line is at the right edge of the scan; a hyphen after "commu" may be cut off
tative) group coefficients.  This is ind\unsure{e|d}ed the basic technical result
enabling them,  from known results on etale cohomology\inserted[typed]{of schemes} (including non
commutative H$^{1}$'s) to deduce corresponding information on homotopy types.
Maybe however \deleted{that} their \hand{observation} has acted rather as a dissuasion for
developping higher non commutative cohomological algebra,  as it seeming-%
\margin{-}ly sa\over{\illegible{1}}{ys} that the non commutative H$^{1}$ ,  plus the commutative H$^{i}$\hand{\unsure{,}}'s, was
all that was needed to recover stringent information about homotopy
types.  In other words,  there was'nt too\xout{\illegible{1}} little in Giraud's book,
but rather,  too much !
